% Options for packages loaded elsewhere
% Options for packages loaded elsewhere
\PassOptionsToPackage{unicode}{hyperref}
\PassOptionsToPackage{hyphens}{url}
\PassOptionsToPackage{dvipsnames,svgnames,x11names}{xcolor}
%
\documentclass[
  letterpaper,
  DIV=11,
  numbers=noendperiod]{scrartcl}
\usepackage{xcolor}
\usepackage{amsmath,amssymb}
\setcounter{secnumdepth}{-\maxdimen} % remove section numbering
\usepackage{iftex}
\ifPDFTeX
  \usepackage[T1]{fontenc}
  \usepackage[utf8]{inputenc}
  \usepackage{textcomp} % provide euro and other symbols
\else % if luatex or xetex
  \usepackage{unicode-math} % this also loads fontspec
  \defaultfontfeatures{Scale=MatchLowercase}
  \defaultfontfeatures[\rmfamily]{Ligatures=TeX,Scale=1}
\fi
\usepackage{lmodern}
\ifPDFTeX\else
  % xetex/luatex font selection
\fi
% Use upquote if available, for straight quotes in verbatim environments
\IfFileExists{upquote.sty}{\usepackage{upquote}}{}
\IfFileExists{microtype.sty}{% use microtype if available
  \usepackage[]{microtype}
  \UseMicrotypeSet[protrusion]{basicmath} % disable protrusion for tt fonts
}{}
\makeatletter
\@ifundefined{KOMAClassName}{% if non-KOMA class
  \IfFileExists{parskip.sty}{%
    \usepackage{parskip}
  }{% else
    \setlength{\parindent}{0pt}
    \setlength{\parskip}{6pt plus 2pt minus 1pt}}
}{% if KOMA class
  \KOMAoptions{parskip=half}}
\makeatother
% Make \paragraph and \subparagraph free-standing
\makeatletter
\ifx\paragraph\undefined\else
  \let\oldparagraph\paragraph
  \renewcommand{\paragraph}{
    \@ifstar
      \xxxParagraphStar
      \xxxParagraphNoStar
  }
  \newcommand{\xxxParagraphStar}[1]{\oldparagraph*{#1}\mbox{}}
  \newcommand{\xxxParagraphNoStar}[1]{\oldparagraph{#1}\mbox{}}
\fi
\ifx\subparagraph\undefined\else
  \let\oldsubparagraph\subparagraph
  \renewcommand{\subparagraph}{
    \@ifstar
      \xxxSubParagraphStar
      \xxxSubParagraphNoStar
  }
  \newcommand{\xxxSubParagraphStar}[1]{\oldsubparagraph*{#1}\mbox{}}
  \newcommand{\xxxSubParagraphNoStar}[1]{\oldsubparagraph{#1}\mbox{}}
\fi
\makeatother

\usepackage{color}
\usepackage{fancyvrb}
\newcommand{\VerbBar}{|}
\newcommand{\VERB}{\Verb[commandchars=\\\{\}]}
\DefineVerbatimEnvironment{Highlighting}{Verbatim}{commandchars=\\\{\}}
% Add ',fontsize=\small' for more characters per line
\usepackage{framed}
\definecolor{shadecolor}{RGB}{241,243,245}
\newenvironment{Shaded}{\begin{snugshade}}{\end{snugshade}}
\newcommand{\AlertTok}[1]{\textcolor[rgb]{0.68,0.00,0.00}{#1}}
\newcommand{\AnnotationTok}[1]{\textcolor[rgb]{0.37,0.37,0.37}{#1}}
\newcommand{\AttributeTok}[1]{\textcolor[rgb]{0.40,0.45,0.13}{#1}}
\newcommand{\BaseNTok}[1]{\textcolor[rgb]{0.68,0.00,0.00}{#1}}
\newcommand{\BuiltInTok}[1]{\textcolor[rgb]{0.00,0.23,0.31}{#1}}
\newcommand{\CharTok}[1]{\textcolor[rgb]{0.13,0.47,0.30}{#1}}
\newcommand{\CommentTok}[1]{\textcolor[rgb]{0.37,0.37,0.37}{#1}}
\newcommand{\CommentVarTok}[1]{\textcolor[rgb]{0.37,0.37,0.37}{\textit{#1}}}
\newcommand{\ConstantTok}[1]{\textcolor[rgb]{0.56,0.35,0.01}{#1}}
\newcommand{\ControlFlowTok}[1]{\textcolor[rgb]{0.00,0.23,0.31}{\textbf{#1}}}
\newcommand{\DataTypeTok}[1]{\textcolor[rgb]{0.68,0.00,0.00}{#1}}
\newcommand{\DecValTok}[1]{\textcolor[rgb]{0.68,0.00,0.00}{#1}}
\newcommand{\DocumentationTok}[1]{\textcolor[rgb]{0.37,0.37,0.37}{\textit{#1}}}
\newcommand{\ErrorTok}[1]{\textcolor[rgb]{0.68,0.00,0.00}{#1}}
\newcommand{\ExtensionTok}[1]{\textcolor[rgb]{0.00,0.23,0.31}{#1}}
\newcommand{\FloatTok}[1]{\textcolor[rgb]{0.68,0.00,0.00}{#1}}
\newcommand{\FunctionTok}[1]{\textcolor[rgb]{0.28,0.35,0.67}{#1}}
\newcommand{\ImportTok}[1]{\textcolor[rgb]{0.00,0.46,0.62}{#1}}
\newcommand{\InformationTok}[1]{\textcolor[rgb]{0.37,0.37,0.37}{#1}}
\newcommand{\KeywordTok}[1]{\textcolor[rgb]{0.00,0.23,0.31}{\textbf{#1}}}
\newcommand{\NormalTok}[1]{\textcolor[rgb]{0.00,0.23,0.31}{#1}}
\newcommand{\OperatorTok}[1]{\textcolor[rgb]{0.37,0.37,0.37}{#1}}
\newcommand{\OtherTok}[1]{\textcolor[rgb]{0.00,0.23,0.31}{#1}}
\newcommand{\PreprocessorTok}[1]{\textcolor[rgb]{0.68,0.00,0.00}{#1}}
\newcommand{\RegionMarkerTok}[1]{\textcolor[rgb]{0.00,0.23,0.31}{#1}}
\newcommand{\SpecialCharTok}[1]{\textcolor[rgb]{0.37,0.37,0.37}{#1}}
\newcommand{\SpecialStringTok}[1]{\textcolor[rgb]{0.13,0.47,0.30}{#1}}
\newcommand{\StringTok}[1]{\textcolor[rgb]{0.13,0.47,0.30}{#1}}
\newcommand{\VariableTok}[1]{\textcolor[rgb]{0.07,0.07,0.07}{#1}}
\newcommand{\VerbatimStringTok}[1]{\textcolor[rgb]{0.13,0.47,0.30}{#1}}
\newcommand{\WarningTok}[1]{\textcolor[rgb]{0.37,0.37,0.37}{\textit{#1}}}

\usepackage{longtable,booktabs,array}
\newcounter{none} % for unnumbered tables
\usepackage{calc} % for calculating minipage widths
% Correct order of tables after \paragraph or \subparagraph
\usepackage{etoolbox}
\makeatletter
\patchcmd\longtable{\par}{\if@noskipsec\mbox{}\fi\par}{}{}
\makeatother
% Allow footnotes in longtable head/foot
\IfFileExists{footnotehyper.sty}{\usepackage{footnotehyper}}{\usepackage{footnote}}
\makesavenoteenv{longtable}
\usepackage{graphicx}
\makeatletter
\newsavebox\pandoc@box
\newcommand*\pandocbounded[1]{% scales image to fit in text height/width
  \sbox\pandoc@box{#1}%
  \Gscale@div\@tempa{\textheight}{\dimexpr\ht\pandoc@box+\dp\pandoc@box\relax}%
  \Gscale@div\@tempb{\linewidth}{\wd\pandoc@box}%
  \ifdim\@tempb\p@<\@tempa\p@\let\@tempa\@tempb\fi% select the smaller of both
  \ifdim\@tempa\p@<\p@\scalebox{\@tempa}{\usebox\pandoc@box}%
  \else\usebox{\pandoc@box}%
  \fi%
}
% Set default figure placement to htbp
\def\fps@figure{htbp}
\makeatother





\setlength{\emergencystretch}{3em} % prevent overfull lines

\providecommand{\tightlist}{%
  \setlength{\itemsep}{0pt}\setlength{\parskip}{0pt}}



 


\KOMAoption{captions}{tableheading}
\makeatletter
\@ifpackageloaded{caption}{}{\usepackage{caption}}
\AtBeginDocument{%
\ifdefined\contentsname
  \renewcommand*\contentsname{Table of contents}
\else
  \newcommand\contentsname{Table of contents}
\fi
\ifdefined\listfigurename
  \renewcommand*\listfigurename{List of Figures}
\else
  \newcommand\listfigurename{List of Figures}
\fi
\ifdefined\listtablename
  \renewcommand*\listtablename{List of Tables}
\else
  \newcommand\listtablename{List of Tables}
\fi
\ifdefined\figurename
  \renewcommand*\figurename{Figure}
\else
  \newcommand\figurename{Figure}
\fi
\ifdefined\tablename
  \renewcommand*\tablename{Table}
\else
  \newcommand\tablename{Table}
\fi
}
\@ifpackageloaded{float}{}{\usepackage{float}}
\floatstyle{ruled}
\@ifundefined{c@chapter}{\newfloat{codelisting}{h}{lop}}{\newfloat{codelisting}{h}{lop}[chapter]}
\floatname{codelisting}{Listing}
\newcommand*\listoflistings{\listof{codelisting}{List of Listings}}
\makeatother
\makeatletter
\makeatother
\makeatletter
\@ifpackageloaded{caption}{}{\usepackage{caption}}
\@ifpackageloaded{subcaption}{}{\usepackage{subcaption}}
\makeatother
\usepackage{bookmark}
\IfFileExists{xurl.sty}{\usepackage{xurl}}{} % add URL line breaks if available
\urlstyle{same}
% fallback for those not using the hyperref driver hyperxmp:
\makeatletter
\@ifundefined{xmpquote}{\newcommand{\xmpquote}[1]{#1}}{}
\makeatother
\hypersetup{
  pdftitle={Predicting Palmer Penguin Sex from Physical Measurements},
  pdfauthor={DSCI 522 Group 36},
  colorlinks=true,
  linkcolor={blue},
  filecolor={Maroon},
  citecolor={Blue},
  urlcolor={Blue},
  pdfcreator={LaTeX via pandoc}}


\title{Predicting Palmer Penguin Sex from Physical Measurements}
\author{DSCI 522 Group 36}
\date{}
\begin{document}
\maketitle

\renewcommand*\contentsname{Table of contents}
{
\setcounter{tocdepth}{3}
\tableofcontents
}

\subsection{Summary}\label{summary}

In this project, we investigate whether the biological sex (male or
female) of adult Palmer penguins can be accurately predicted from four
measurements: bill length, bill depth, flipper length, and body mass.
Using data from 333 sexed penguins across three species (Adélie,
Chinstrap, and Gentoo), we trained a separate k-Nearest Neighbors (k-NN)
classifier for each species, choosing the number of neighbors with
5-fold cross-validation. On the held-out test sets, our models correctly
classified 32 of 37 Adélie penguins (86.5\%), 16 of 17 Chinstrap
penguins (94.1\%), and 29 of 30 Gentoo penguins (96.7\%), well above the
roughly 50\% accuracy of always predicting the more common sex. These
results show that sexual dimorphism in body size is strong enough to
classify penguin sex reliably without genetic testing. However, because
some males and females overlap in size and our test sets are small,
particularly for Chinstrap penguins, we recommend using the model as a
quick first estimate in the field, with DNA testing reserved for birds
whose predicted sex is uncertain.

\subsection{Introduction}\label{introduction}

Determining the biological sex of seabirds is essential for ecological
monitoring and population modeling, as males and females often exhibit
different behaviors, survival rates, and energetic demands during
breeding seasons (Williams, 1995; Gorman et al., 2014). However, many
penguin species are considered sexually monomorphic to the human eye,
meaning males and females share identical plumage patterns and cannot be
reliably distinguished visually in the field (Gorman et al., 2014).

Currently, the most reliable way to determine a penguin's sex is through
molecular testing, which requires collecting blood or feather samples,
storing them in remote field conditions, and sending them to a
laboratory for analysis. This process can take weeks, adds cost, and
requires additional handling that can stress the birds. Observing
breeding behavior can also reveal sex, but only during certain parts of
the breeding season and only after long periods of observation. As a
result, researchers often have to wait or rely on incomplete information
when studying how males and females differ in their foraging, survival,
or response to environmental change. This matters especially in
Antarctica, where rapid warming along the Antarctic Peninsula may affect
males and females differently, and where accurate sex ratios are needed
to judge the health of breeding populations. If external body
measurements can accurately predict sex, field researchers can rapidly
classify individuals using simple, non-invasive calipers and scales. In
this report, we ask: \textbf{Can we accurately predict the sex (male or
female) of Adélie, Chinstrap, and Gentoo penguins from their bill
length, bill depth, flipper length, and body mass using species-specific
k-Nearest Neighbors (k-NN) classification models?}

We use the Palmer Archipelago (Antarctica) penguin dataset originally
collected by Dr.~Kristen Gorman and the Palmer Station Long Term
Ecological Research (LTER) program (Gorman et al., 2014), and curated
for data science by Horst et al.~(2020). Each row in the dataset
represents an individual adult penguin observed on one of three islands
(Biscoe, Dream, or Torgersen) between 2007 and 2009. From the available
variables, we use species (Adélie, Chinstrap, or Gentoo) to stratify our
models, four continuous physical predictors (bill length and bill depth
in millimeters, flipper length in millimeters, and body mass in grams),
and our binary target variable, biological sex, which takes the values
female or male. This dataset is well suited to our question because sex
was recorded for almost every penguin, and all four measurements were
taken on the same individuals. Of the 344 penguins in the dataset, 333
have both a recorded sex and complete measurements, giving us labeled
examples from which a model can learn the relationship between body size
and sex.

\subsection{Methods}\label{methods}

\subsubsection{Data}\label{data}

The data were collected by Dr.~Kristen Gorman and the Palmer Station
Long Term Ecological Research (LTER) program between 2007 and 2009
(Gorman et al., 2014), and curated for data science as the
palmerpenguins data set (Horst et al., 2020), which can be found
\href{https://allisonhorst.github.io/palmerpenguins/}{here}. To ensure
full reproducibility, our analysis downloads the data directly from the
palmerpenguins repository and saves an unmodified copy to
\texttt{data/raw/}. Observations without a recorded sex cannot be used
for supervised training or evaluation, so we remove them, along with the
two penguins that are missing all body measurements. This leaves 333
penguins, and we save the cleaned data to \texttt{data/processed/}.

\begin{Shaded}
\begin{Highlighting}[]
\ImportTok{import}\NormalTok{ os}

\ImportTok{import}\NormalTok{ matplotlib.pyplot }\ImportTok{as}\NormalTok{ plt}
\ImportTok{import}\NormalTok{ pandas }\ImportTok{as}\NormalTok{ pd}
\ImportTok{import}\NormalTok{ requests}
\ImportTok{import}\NormalTok{ seaborn }\ImportTok{as}\NormalTok{ sns}
\ImportTok{from}\NormalTok{ sklearn.impute }\ImportTok{import}\NormalTok{ SimpleImputer}
\ImportTok{from}\NormalTok{ sklearn.metrics }\ImportTok{import}\NormalTok{ ConfusionMatrixDisplay, classification\_report}
\ImportTok{from}\NormalTok{ sklearn.model\_selection }\ImportTok{import}\NormalTok{ GridSearchCV, train\_test\_split}
\ImportTok{from}\NormalTok{ sklearn.neighbors }\ImportTok{import}\NormalTok{ KNeighborsClassifier}
\ImportTok{from}\NormalTok{ sklearn.pipeline }\ImportTok{import}\NormalTok{ make\_pipeline}
\ImportTok{from}\NormalTok{ sklearn.preprocessing }\ImportTok{import}\NormalTok{ StandardScaler}

\CommentTok{\# create the data folders if they don\textquotesingle{}t exist}
\CommentTok{\# (git does not track empty folders)}
\NormalTok{os.makedirs(}\StringTok{"../data/raw"}\NormalTok{, exist\_ok}\OperatorTok{=}\VariableTok{True}\NormalTok{)}
\NormalTok{os.makedirs(}\StringTok{"../data/processed"}\NormalTok{, exist\_ok}\OperatorTok{=}\VariableTok{True}\NormalTok{)}

\CommentTok{\# download the Palmer penguins data (CSV) from the palmerpenguins repository}
\NormalTok{url }\OperatorTok{=} \StringTok{"https://raw.githubusercontent.com/allisonhorst/palmerpenguins/main/inst/extdata/penguins.csv"}

\NormalTok{request }\OperatorTok{=}\NormalTok{ requests.get(url)}
\ControlFlowTok{with} \BuiltInTok{open}\NormalTok{(}\StringTok{"../data/raw/penguins.csv"}\NormalTok{, }\StringTok{"wb"}\NormalTok{) }\ImportTok{as}\NormalTok{ f:}
\NormalTok{    f.write(request.content)}

\NormalTok{df\_raw }\OperatorTok{=}\NormalTok{ pd.read\_csv(}\StringTok{"../data/raw/penguins.csv"}\NormalTok{)}

\CommentTok{\# keep only the variables used in the analysis and drop rows with missing}
\CommentTok{\# values (penguins with no recorded sex cannot be used to train or test the model)}
\NormalTok{cols\_needed }\OperatorTok{=}\NormalTok{ [}\StringTok{"species"}\NormalTok{, }\StringTok{"sex"}\NormalTok{, }\StringTok{"flipper\_length\_mm"}\NormalTok{, }\StringTok{"bill\_length\_mm"}\NormalTok{,}
               \StringTok{"bill\_depth\_mm"}\NormalTok{, }\StringTok{"body\_mass\_g"}\NormalTok{]}
\NormalTok{df\_clean }\OperatorTok{=}\NormalTok{ df\_raw[cols\_needed].dropna().reset\_index(drop}\OperatorTok{=}\VariableTok{True}\NormalTok{)}
\NormalTok{df\_clean.to\_csv(}\StringTok{"../data/processed/penguins\_clean.csv"}\NormalTok{, index}\OperatorTok{=}\VariableTok{False}\NormalTok{)}
\end{Highlighting}
\end{Shaded}

\subsubsection{Analysis}\label{analysis}

To uphold the golden rule of machine learning---ensuring that test
observations never influence exploratory decisions, feature scaling, or
hyperparameter selection--- we split the cleaned dataset into training
(75\%) and test (25\%) subsets prior to performing any exploratory data
analysis or preprocessing. Because body size differs substantially
across penguin species (for instance, a female Gentoo penguin is
typically heavier than a male Adelie penguin), we partition the data
separately for each species using stratified random sampling to preserve
equal proportions of males and females in both sets. The split data is
saved in the processed folder as well with the cleaned data.

\begin{Shaded}
\begin{Highlighting}[]
\CommentTok{\# split each species before exploratory analysis}
\NormalTok{features }\OperatorTok{=}\NormalTok{ [}\StringTok{"flipper\_length\_mm"}\NormalTok{, }\StringTok{"bill\_length\_mm"}\NormalTok{, }\StringTok{"bill\_depth\_mm"}\NormalTok{, }\StringTok{"body\_mass\_g"}\NormalTok{]}
\NormalTok{target }\OperatorTok{=} \StringTok{"sex"}
\NormalTok{species\_list }\OperatorTok{=}\NormalTok{ [}\StringTok{"Adelie"}\NormalTok{, }\StringTok{"Chinstrap"}\NormalTok{, }\StringTok{"Gentoo"}\NormalTok{]}

\NormalTok{splits }\OperatorTok{=}\NormalTok{ \{\}}
\NormalTok{train\_frames, test\_frames }\OperatorTok{=}\NormalTok{ [], []}
\ControlFlowTok{for}\NormalTok{ sp }\KeywordTok{in}\NormalTok{ species\_list:}
\NormalTok{    sub }\OperatorTok{=}\NormalTok{ df\_clean[df\_clean[}\StringTok{"species"}\NormalTok{] }\OperatorTok{==}\NormalTok{ sp]}
\NormalTok{    X\_train, X\_test, y\_train, y\_test }\OperatorTok{=}\NormalTok{ train\_test\_split(}
\NormalTok{        sub[features], sub[target], test\_size}\OperatorTok{=}\FloatTok{0.25}\NormalTok{, stratify}\OperatorTok{=}\NormalTok{sub[target], random\_state}\OperatorTok{=}\DecValTok{522}
\NormalTok{    )}
\NormalTok{    splits[sp] }\OperatorTok{=}\NormalTok{ (X\_train, X\_test, y\_train, y\_test)}
\NormalTok{    train\_frames.append(X\_train.assign(sex}\OperatorTok{=}\NormalTok{y\_train, species}\OperatorTok{=}\NormalTok{sp))}
\NormalTok{    test\_frames.append(X\_test.assign(sex}\OperatorTok{=}\NormalTok{y\_test, species}\OperatorTok{=}\NormalTok{sp))}

\NormalTok{train\_eda }\OperatorTok{=}\NormalTok{ pd.concat(train\_frames, ignore\_index}\OperatorTok{=}\VariableTok{True}\NormalTok{)}
\NormalTok{test\_df }\OperatorTok{=}\NormalTok{ pd.concat(test\_frames, ignore\_index}\OperatorTok{=}\VariableTok{True}\NormalTok{)}
\NormalTok{train\_eda.to\_csv(}\StringTok{"../data/processed/penguins\_train.csv"}\NormalTok{, index}\OperatorTok{=}\VariableTok{False}\NormalTok{)}
\NormalTok{test\_df.to\_csv(}\StringTok{"../data/processed/penguins\_test.csv"}\NormalTok{, index}\OperatorTok{=}\VariableTok{False}\NormalTok{)}
\end{Highlighting}
\end{Shaded}

We used a k-Nearest Neighbors (k-NN) classifier for each species (Cover
\& Hart, 1967), which predicts a penguin's sex as the most common sex
among the k training penguins with the most similar measurements. This
suits our question because males and females form fairly distinct groups
when their measurements are combined, and k-NN is simple to explain to
field researchers. Because k-NN relies on distances, all four
measurements were standardized (with median imputation as a safeguard
against missing values) inside a pipeline, so these steps were learned
from the training data only. We chose k using 5-fold cross-validation on
the training set, trying odd values from 1 to 19, and used the test set
only once, to evaluate the final models.

We evaluated the models using accuracy, the proportion of penguins whose
sex is predicted correctly. Accuracy is appropriate because males and
females are roughly balanced within each species and both types of error
are equally costly; we compare it with a baseline of about 50\%, the
accuracy of always predicting the more common sex, and also report
precision and recall for each sex. The analysis was performed in Python
(Van Rossum \& Drake, 2009) using requests (Reitz, 2011), pandas
(McKinney, 2010), scikit-learn (Pedregosa et al., 2011), Matplotlib
(Hunter, 2007), and seaborn (Waskom, 2021); the code is available at
https://github.com/UBC-MDS-2026-27/DSCI-522-Group-36-Palmer-Penguins.

\subsection{Results}\label{results}

\subsubsection{Exploratory Analysis}\label{exploratory-analysis}

Before building our models, we explored the training data to check
whether male and female penguins differ in the four physical
measurements. This step is essential to our question since if males and
females had similar measurements, there would be little reason to expect
a classifier to predict sex from them. All exploratory analysis uses the
training set only, so the test set does not influence any of our
decisions.

Table~\ref{tbl-summary} shows the number of penguins and the mean of
each measurement for each species and sex in the training set.

\begin{Shaded}
\begin{Highlighting}[]
\NormalTok{summary\_table }\OperatorTok{=}\NormalTok{ (}
\NormalTok{    train\_eda.groupby([}\StringTok{"species"}\NormalTok{, }\StringTok{"sex"}\NormalTok{])}
\NormalTok{    .agg(}
\NormalTok{        count}\OperatorTok{=}\NormalTok{(}\StringTok{"body\_mass\_g"}\NormalTok{, }\StringTok{"size"}\NormalTok{),}
\NormalTok{        body\_mass}\OperatorTok{=}\NormalTok{(}\StringTok{"body\_mass\_g"}\NormalTok{, }\StringTok{"mean"}\NormalTok{),}
\NormalTok{        flipper\_length}\OperatorTok{=}\NormalTok{(}\StringTok{"flipper\_length\_mm"}\NormalTok{, }\StringTok{"mean"}\NormalTok{),}
\NormalTok{        bill\_length}\OperatorTok{=}\NormalTok{(}\StringTok{"bill\_length\_mm"}\NormalTok{, }\StringTok{"mean"}\NormalTok{),}
\NormalTok{        bill\_depth}\OperatorTok{=}\NormalTok{(}\StringTok{"bill\_depth\_mm"}\NormalTok{, }\StringTok{"mean"}\NormalTok{),}
\NormalTok{    )}
\NormalTok{    .}\BuiltInTok{round}\NormalTok{(}\DecValTok{1}\NormalTok{)}
\NormalTok{    .rename(columns}\OperatorTok{=}\NormalTok{\{}
        \StringTok{"count"}\NormalTok{: }\StringTok{"Count"}\NormalTok{,}
        \StringTok{"body\_mass"}\NormalTok{: }\StringTok{"Body mass (g)"}\NormalTok{,}
        \StringTok{"flipper\_length"}\NormalTok{: }\StringTok{"Flipper length (mm)"}\NormalTok{,}
        \StringTok{"bill\_length"}\NormalTok{: }\StringTok{"Bill length (mm)"}\NormalTok{,}
        \StringTok{"bill\_depth"}\NormalTok{: }\StringTok{"Bill depth (mm)"}\NormalTok{,}
\NormalTok{    \})}
\NormalTok{)}
\NormalTok{summary\_table.index.names }\OperatorTok{=}\NormalTok{ [}\StringTok{"Species"}\NormalTok{, }\StringTok{"Sex"}\NormalTok{]}
\NormalTok{sex\_order }\OperatorTok{=} \BuiltInTok{sorted}\NormalTok{(train\_eda[}\StringTok{"sex"}\NormalTok{].unique())}
\NormalTok{summary\_table}
\end{Highlighting}
\end{Shaded}

{

\begin{longtable}[]{@{}lllllll@{}}

\caption{\label{tbl-summary}Number of penguins and mean physical
measurements by species and sex (training set).}

\tabularnewline

\toprule\noalign{}
& & Count & Body mass (g) & Flipper length (mm) & Bill length (mm) &
Bill depth (mm) \\
Species & Sex & & & & & \\
\midrule\noalign{}
\endhead
\bottomrule\noalign{}
\endlastfoot
\multirow{2}{*}{Adelie} & female & 54 & 3368.5 & 188.2 & 37.4 & 17.6 \\
& male & 55 & 4056.4 & 192.1 & 40.2 & 19.0 \\
\multirow{2}{*}{Chinstrap} & female & 25 & 3550.0 & 192.0 & 46.8 &
17.7 \\
& male & 26 & 3930.8 & 199.3 & 50.9 & 19.2 \\
\multirow{2}{*}{Gentoo} & female & 43 & 4689.0 & 212.3 & 45.6 & 14.2 \\
& male & 46 & 5454.9 & 222.0 & 49.6 & 15.8 \\

\end{longtable}

}

Within each species, males are on average heavier than females and tend
to have longer flippers, longer bills, and deeper bills. The numbers of
males and females are close to equal in every species, so a model that
always guessed the same sex would be correct only about half of the
time; this gives a baseline of roughly 50\% accuracy for our models to
beat. The table also shows that Gentoo penguins are much larger than
Adélie and Chinstrap penguins, which supports our decision to build a
separate model for each species.

Figure~\ref{fig-mass-box} shows the distribution of body mass by sex
within each species. In every species, the typical male is heavier than
the typical female. However, the boxes overlap, meaning that some large
females are heavier than some small males. Body mass alone therefore
cannot perfectly separate the sexes.

\begin{Shaded}
\begin{Highlighting}[]
\NormalTok{fig, ax }\OperatorTok{=}\NormalTok{ plt.subplots(figsize}\OperatorTok{=}\NormalTok{(}\DecValTok{8}\NormalTok{, }\DecValTok{4}\NormalTok{))}
\NormalTok{sns.boxplot(}
\NormalTok{    data}\OperatorTok{=}\NormalTok{train\_eda, x}\OperatorTok{=}\StringTok{"species"}\NormalTok{, y}\OperatorTok{=}\StringTok{"body\_mass\_g"}\NormalTok{, hue}\OperatorTok{=}\StringTok{"sex"}\NormalTok{,}
\NormalTok{    order}\OperatorTok{=}\NormalTok{species\_list, hue\_order}\OperatorTok{=}\NormalTok{sex\_order, ax}\OperatorTok{=}\NormalTok{ax,}
\NormalTok{)}
\NormalTok{ax.set\_xlabel(}\StringTok{"Species"}\NormalTok{)}
\NormalTok{ax.set\_ylabel(}\StringTok{"Body mass (g)"}\NormalTok{)}
\NormalTok{ax.legend(title}\OperatorTok{=}\StringTok{"Sex"}\NormalTok{)}
\NormalTok{plt.tight\_layout()}
\NormalTok{plt.show()}
\end{Highlighting}
\end{Shaded}

\begin{figure}[H]

\centering{

\pandocbounded{\includegraphics[keepaspectratio]{penguins-analysis_files/figure-pdf/fig-mass-box-output-1.pdf}}

}

\caption{\label{fig-mass-box}Body mass of female and male penguins in
each species (training set). Males are heavier on average, but the
ranges overlap.}

\end{figure}%

Figure~\ref{fig-scatter} shows that combining two measurements separates
the sexes more clearly than one measurement on its own. When bill depth
is plotted against body mass, males and females form fairly distinct
groups within each species. This suggests that a model using all four
measurements together should be able to predict sex well. It also
motivates our choice of k-NN, which classifies a penguin based on the
sex of the most similar penguins in the training data: if males and
females form separate groups, a penguin's nearest neighbors are likely
to share its sex.

\begin{Shaded}
\begin{Highlighting}[]
\NormalTok{g }\OperatorTok{=}\NormalTok{ sns.relplot(}
\NormalTok{    data}\OperatorTok{=}\NormalTok{train\_eda, x}\OperatorTok{=}\StringTok{"body\_mass\_g"}\NormalTok{, y}\OperatorTok{=}\StringTok{"bill\_depth\_mm"}\NormalTok{, hue}\OperatorTok{=}\StringTok{"sex"}\NormalTok{,}
\NormalTok{    col}\OperatorTok{=}\StringTok{"species"}\NormalTok{, col\_order}\OperatorTok{=}\NormalTok{species\_list, hue\_order}\OperatorTok{=}\NormalTok{sex\_order,}
\NormalTok{    height}\OperatorTok{=}\FloatTok{3.2}\NormalTok{, aspect}\OperatorTok{=}\FloatTok{0.9}\NormalTok{, alpha}\OperatorTok{=}\FloatTok{0.7}\NormalTok{,}
\NormalTok{    facet\_kws}\OperatorTok{=}\NormalTok{\{}\StringTok{"sharex"}\NormalTok{: }\VariableTok{False}\NormalTok{, }\StringTok{"sharey"}\NormalTok{: }\VariableTok{False}\NormalTok{\},}
\NormalTok{)}
\NormalTok{g.set\_axis\_labels(}\StringTok{"Body mass (g)"}\NormalTok{, }\StringTok{"Bill depth (mm)"}\NormalTok{)}
\NormalTok{g.set\_titles(}\StringTok{"}\SpecialCharTok{\{col\_name\}}\StringTok{"}\NormalTok{)}
\NormalTok{g.legend.set\_title(}\StringTok{"Sex"}\NormalTok{)}
\NormalTok{plt.show()}
\end{Highlighting}
\end{Shaded}

\begin{figure}[H]

\centering{

\pandocbounded{\includegraphics[keepaspectratio]{penguins-analysis_files/figure-pdf/fig-scatter-output-1.pdf}}

}

\caption{\label{fig-scatter}Bill depth against body mass for each
species, colored by sex (training set). Males and females form fairly
distinct groups when two measurements are combined.}

\end{figure}%

Finally, because body mass is measured in grams (values in the
thousands) while the other measurements are in millimeters (values in
the tens or hundreds), the features are on very different scales. Since
k-NN classifies penguins using distances between their measurements,
body mass would dominate those distances if left unscaled. This is why
our modeling pipeline standardizes all four measurements before fitting
the classifier.

\subsubsection{Model}\label{model}

\begin{Shaded}
\begin{Highlighting}[]
\CommentTok{\#Question: Can we predict a penguin\textquotesingle{}s sex from its measurements (flipper length, bill length, bill depth, body mass) using k{-}NN}

\NormalTok{pipe }\OperatorTok{=}\NormalTok{ make\_pipeline(}
\NormalTok{    SimpleImputer(strategy}\OperatorTok{=}\StringTok{"median"}\NormalTok{),}
\NormalTok{    StandardScaler(),}
\NormalTok{    KNeighborsClassifier(),}
\NormalTok{)}
\NormalTok{param\_grid }\OperatorTok{=}\NormalTok{ \{}\StringTok{"kneighborsclassifier\_\_n\_neighbors"}\NormalTok{: }\BuiltInTok{range}\NormalTok{(}\DecValTok{1}\NormalTok{, }\DecValTok{20}\NormalTok{, }\DecValTok{2}\NormalTok{)\}}

\NormalTok{results }\OperatorTok{=}\NormalTok{ []}

\ControlFlowTok{for}\NormalTok{ species }\KeywordTok{in}\NormalTok{ [}\StringTok{"Adelie"}\NormalTok{, }\StringTok{"Chinstrap"}\NormalTok{, }\StringTok{"Gentoo"}\NormalTok{]:}
    \CommentTok{\#reuse the training/test split made in Methods}
\NormalTok{    X\_train, X\_test, y\_train, y\_test }\OperatorTok{=}\NormalTok{ splits[species]}

\NormalTok{    grid }\OperatorTok{=}\NormalTok{ GridSearchCV(pipe, param\_grid}\OperatorTok{=}\NormalTok{param\_grid, cv}\OperatorTok{=}\DecValTok{5}\NormalTok{, scoring}\OperatorTok{=}\StringTok{"accuracy"}\NormalTok{)}
\NormalTok{    grid.fit(X\_train, y\_train)}

\NormalTok{    best\_k }\OperatorTok{=}\NormalTok{ grid.best\_params\_[}\StringTok{"kneighborsclassifier\_\_n\_neighbors"}\NormalTok{]}
\NormalTok{    test\_acc }\OperatorTok{=}\NormalTok{ grid.score(X\_test, y\_test)}
\NormalTok{    y\_pred }\OperatorTok{=}\NormalTok{ grid.predict(X\_test)}

\NormalTok{    results.append(\{}\StringTok{"Species"}\NormalTok{: species, }\StringTok{"Best k"}\NormalTok{: best\_k, }\StringTok{"CV Accuracy"}\NormalTok{: }\BuiltInTok{round}\NormalTok{(grid.best\_score\_, }\DecValTok{3}\NormalTok{), }\StringTok{"Test Accuracy"}\NormalTok{: }\BuiltInTok{round}\NormalTok{(test\_acc, }\DecValTok{3}\NormalTok{),}
                    \StringTok{"Correct"}\NormalTok{: }\BuiltInTok{int}\NormalTok{((y\_pred }\OperatorTok{==}\NormalTok{ y\_test).}\BuiltInTok{sum}\NormalTok{()), }\StringTok{"Test size"}\NormalTok{: }\BuiltInTok{len}\NormalTok{(y\_test)\})}

    \BuiltInTok{print}\NormalTok{(}\SpecialStringTok{f"=== }\SpecialCharTok{\{}\NormalTok{species}\SpecialCharTok{\}}\SpecialStringTok{ (Best k=}\SpecialCharTok{\{}\NormalTok{best\_k}\SpecialCharTok{\}}\SpecialStringTok{) ==="}\NormalTok{)}
    \BuiltInTok{print}\NormalTok{(classification\_report(y\_test, y\_pred))}
\end{Highlighting}
\end{Shaded}

\begin{verbatim}
=== Adelie (Best k=9) ===
              precision    recall  f1-score   support

      female       0.85      0.89      0.87        19
        male       0.88      0.83      0.86        18

    accuracy                           0.86        37
   macro avg       0.87      0.86      0.86        37
weighted avg       0.87      0.86      0.86        37

=== Chinstrap (Best k=9) ===
              precision    recall  f1-score   support

      female       1.00      0.89      0.94         9
        male       0.89      1.00      0.94         8

    accuracy                           0.94        17
   macro avg       0.94      0.94      0.94        17
weighted avg       0.95      0.94      0.94        17

=== Gentoo (Best k=9) ===
              precision    recall  f1-score   support

      female       0.94      1.00      0.97        15
        male       1.00      0.93      0.97        15

    accuracy                           0.97        30
   macro avg       0.97      0.97      0.97        30
weighted avg       0.97      0.97      0.97        30
\end{verbatim}

\begin{Shaded}
\begin{Highlighting}[]
\CommentTok{\# Summary comparison table across all 3 species}
\NormalTok{summary\_df }\OperatorTok{=}\NormalTok{ pd.DataFrame(results)}
\NormalTok{summary\_df}
\end{Highlighting}
\end{Shaded}

{

\begin{longtable}[]{@{}lllllll@{}}

\caption{\label{tbl-results}Selected number of neighbors,
cross-validation (CV) accuracy, and test accuracy for each species.}

\tabularnewline

\toprule\noalign{}
& Species & Best k & CV Accuracy & Test Accuracy & Correct & Test
size \\
\midrule\noalign{}
\endhead
\bottomrule\noalign{}
\endlastfoot
0 & Adelie & 9 & 0.945 & 0.865 & 32 & 37 \\
1 & Chinstrap & 9 & 0.900 & 0.941 & 16 & 17 \\
2 & Gentoo & 9 & 0.944 & 0.967 & 29 & 30 \\

\end{longtable}

}

Cross-validation selected k = 9 for all three species. As summarized in
Table~\ref{tbl-results}, the final models correctly classified 32 of 37
Adélie penguins (86.5\%), 16 of 17 Chinstrap penguins (94.1\%), and 29
of 30 Gentoo penguins (96.7\%) in the test set, well above the roughly
50\% accuracy of always predicting the more common sex. Test accuracy
was close to or slightly above cross-validation accuracy for Chinstrap
and Gentoo penguins, but lower for Adélie penguins (86.5\% compared with
94.5\%), which may reflect the variability of a small test set. The
precision and recall values are {[}similar / different{]} for the two
sexes, which indicates that the models {[}do not / do{]} systematically
favor one sex.

\subsection{Discussion}\label{discussion}

Our results confirm that external physical measurements can accurately
predict the biological sex of adult Palmer penguins when stratified by
species, achieving high held-out test accuracies for all penguins. Our
findings directly answer our research question by demonstrating that
bill length, bill depth, flipper length, and body mass carry sufficient
sex-specific signal to classify male and female penguins reliably. These
results are consistent with our exploratory analysis, which showed that
males are larger than females on average in every species
(Table~\ref{tbl-summary}) and that the sexes form fairly distinct groups
when measurements are combined (Figure~\ref{fig-scatter}). The remaining
errors are also expected: Figure~\ref{fig-mass-box} showed that the body
sizes of males and females overlap, so a small male or a large female
can be misclassified however well the model is tuned.

Despite strong performance, several methodological assumptions and
limitations must be considered. First, our workflow assumes that a
penguin's species is already known before predicting its sex; while
visual species identification is straightforward in the field due to
distinct head markings, any species misidentification would degrade sex
classification accuracy. Second, stratifying the dataset by species
reduced our sample sizes, leaving only 17 Chinstrap penguins in the test
set, where a single misclassification shifts test accuracy by nearly six
percentage points. Third, because k-Nearest Neighbors relies on
Euclidean distances across physical measurements, seasonal weight
fluctuations, such as mass loss during fasting and incubation, or
regional size differences could shift feature distributions and reduce
accuracy if the model is applied outside the Palmer Archipelago without
local recalibration.

The broader impact of this approach lies in making wildlife demographic
monitoring faster, cheaper, and less invasive. Accurate sex ratios are
critical for estimating breeding population viability and detecting
sex-biased mortality from climate change or fisheries interactions.

For future work, it would be interesting to test if we can still predict
sex accurately without using body mass, since weight changes so much
across seasons---for example, by using bill shape ratios or island
location instead. Researchers could also compare k-Nearest Neighbors
with models like logistic regression to get probability scores, which
would help biologists see when a prediction is uncertain and needs a DNA
test to be sure.

\subsection{References}\label{references}

Cover, T., \& Hart, P. (1967). Nearest neighbor pattern classification.
\emph{IEEE Transactions on Information Theory}, 13(1), 21--27.
https://doi.org/10.1109/TIT.1967.1053964

Gorman, K. B., Williams, T. D., \& Fraser, W. R. (2014). Ecological
sexual dimorphism and environmental variability within a community of
Antarctic penguins (genus \emph{Pygoscelis}). \emph{PLoS ONE}, 9(3),
e90081. https://doi.org/10.1371/journal.pone.0090081

Horst, A. M., Hill, A. P., \& Gorman, K. B. (2020).
\emph{palmerpenguins: Palmer Archipelago (Antarctica) penguin data} (R
package version 0.1.0). https://allisonhorst.github.io/palmerpenguins/

Hunter, J. D. (2007). Matplotlib: A 2D graphics environment.
\emph{Computing in Science \& Engineering}, 9(3), 90--95.
https://doi.org/10.1109/MCSE.2007.55

McKinney, W. (2010). Data structures for statistical computing in
Python. In \emph{Proceedings of the 9th Python in Science Conference}
(pp.~56--61). https://doi.org/10.25080/Majora-92bf1922-00a

Pedregosa, F., Varoquaux, G., Gramfort, A., Michel, V., Thirion, B.,
Grisel, O., Blondel, M., Prettenhofer, P., Weiss, R., Dubourg, V.,
Vanderplas, J., Passos, A., Cournapeau, D., Brucher, M., Perrot, M., \&
Duchesnay, É. (2011). Scikit-learn: Machine learning in Python.
\emph{Journal of Machine Learning Research}, 12, 2825--2830.

Reitz, K. (2011). Requests: HTTP for humans.
https://requests.readthedocs.io/

Van Rossum, G., \& Drake, F. L. (2009). \emph{Python 3 reference
manual}. CreateSpace.

Waskom, M. L. (2021). seaborn: Statistical data visualization.
\emph{Journal of Open Source Software}, 6(60), 3021.
https://doi.org/10.21105/joss.03021

Williams, T. D. (1995). \emph{The penguins: Spheniscidae}. Oxford
University Press.




\end{document}
